% 题干转写：完成
% 解析转写：完成
% 图形核对：完成
\section{第39届复赛第4题：旋转圆环上铰接细杆的平衡与稳定性}

\subsection{题干}
如图\ref{fig:cpho39q04structure}，长度为 $2a$、质量为 $m$ 的匀质刚性细杆的一端以小孔套在半径为 $R$ 的水平圆环固定点 $P$ 上。圆环及 $P$ 点以角速度 $\omega$ 绕过圆心 $O$ 的竖直轴转动；细杆可绕 $P$ 无摩擦转动，并始终位于 $OP$ 所在竖直平面内。细杆与竖直方向夹角为 $\theta$，重力加速度为 $g$。

\begin{figure}[htbp]\centering
  \begin{tikzpicture}[>=Stealth,scale=0.9]
  \draw[thick] (0,0) ellipse (2.25 and 0.58);
  \draw[densely dashed,thick] (0,-0.9) -- (0,1.52);
  \node[below] at (0,-0.12) {$O$};
  \coordinate (P) at (2.15,0);
  \fill (P) circle (1.5pt) node[below right] {$P$};
  \draw[thick] (0,0) -- (P) node[midway,below] {$R$};
  \draw[very thick,fill=gray!18] (P) -- ++(38:2.35) coordinate (B);
  \draw[<->,thin] ($(P)+(38:0.18)$) -- ($(B)+(38:0.18)$) node[midway,above] {$2a$};
  \draw[densely dashed] (P) -- ++(0,1.55);
  \draw[->] ($(P)+(0,0.72)$) arc[start angle=90,end angle=52,radius=0.72]
    node[midway,right] {$\theta$};
  \coordinate (C) at ($(P)!0.5!(B)$);
  \draw[->,thick] (C) -- ++(0,-1.18) node[right] {$mg$};
  \draw[->,thick] (-0.15,0.9) arc[start angle=125,end angle=55,radius=1.18]
    node[midway,above] {$\omega$};
  \node[below=1.05cm of P] {原题图4a：水平圆环与铰接细杆};
\end{tikzpicture}

  \caption{第39届复赛第4题原题图4a的TikZ重绘：转动圆环上的铰接细杆。}\label{fig:cpho39q04structure}
\end{figure}

求：（1）分别在实验室系 $L$ 与随 $OP$ 转动的参考系 $S$ 中写出细杆机械能；（2）在 $S$ 系导出平衡条件；（3）图解分析四个象限中平衡位形数及参数条件；（4）分象限画出对 $P$ 点力矩有贡献的受力示意并检验图解结论；（5）讨论各平衡位形的稳定性。

\subsection{参考解答整理}
\paragraph{（1）两参考系的机械能。}
细杆质心到转轴的距离为 $R+a\sin\theta$。以圆环平面为重力势能零点，实验室系机械能可写为
\[
E_L=mga\cos\theta+\frac23ma^2\dot\theta^2
+\frac12m\left[(R+a\sin\theta)^2+\frac13a^2\sin^2\theta\right]\omega^2.
\]
在随 $OP$ 转动的 $S$ 系中，杆绕端点的转动惯量为 $I_P=4ma^2/3$；还应计入离心势能。因此
\[
\boxed{E_S=\frac23ma^2\dot\theta^2+mga\cos\theta
-\frac12m\left[(R+a\sin\theta)^2+\frac13a^2\sin^2\theta\right]\omega^2.}
\]

\paragraph{（2）平衡条件。}
令有效势能对 $\theta$ 的导数为零，或直接令相对 $P$ 的合力矩为零，均得
\[
mga\sin\theta+ma\omega^2\left(R+\frac43a\sin\theta\right)\cos\theta=0,
\]
即
\[
\boxed{\tan\theta+\frac{4a\omega^2}{3g}\sin\theta+\frac{R\omega^2}{g}=0.}
\]

\paragraph{（3）图解判定。}
令 $A=4a\omega^2/(3g)$、$B=R\omega^2/g$，则 $A,B>0$，平衡方程化为
\[-\tan\theta=A\sin\theta+B.\]
图\ref{fig:cpho39q04tan}重绘答案的图解曲线。另将原方程写为
\[-\cot\theta=A'\cos\theta+B',\qquad A'=\frac{4a}{3R},\quad B'=\frac{g}{R\omega^2},\]
也得到图\ref{fig:cpho39q04cot}的等价图解。
\begin{figure}[htbp]\centering
  \begin{tikzpicture}[>=Stealth,x=1.32cm,y=1.04cm]
  \draw[->] (-0.28,0) -- (5.1,0) node[right] {$\theta$};
  \draw[->] (0,-2.35) -- (0,2.45) node[above] {$y$};
  \foreach \x/\lab in {1.57/{\pi/2},3.14/{\pi},4.71/{3\pi/2}} {
    \draw[densely dashed,gray] (\x,-2.25) -- (\x,2.3);
    \node[below] at (\x,0) {$\lab$};
  }
  \begin{scope}\clip (0,-2.25) rectangle (5.0,2.3);
    \draw[thick,domain=0.06:1.50,samples=100] plot (\x,{-tan(\x r)});
    \draw[thick,domain=1.64:4.62,samples=150] plot (\x,{-tan(\x r)});
    \draw[thick,domain=4.79:4.98,samples=40] plot (\x,{-tan(\x r)});
    \draw[thick,blue!65!black,domain=0:5.0,samples=150] plot (\x,{0.88*sin(\x r)+0.72});
  \end{scope}
  \draw[densely dotted] (0,1.60) -- (5.0,1.60) node[right] {$A+B$};
  \draw[densely dotted] (0,-0.16) -- (5.0,-0.16) node[right] {$B-A$};
  \node[blue!65!black,above] at (3.85,0.17) {$A\sin\theta+B$};
  \node[above] at (4.35,1.68) {$-\tan\theta$};
  \node[below=0.68cm] at (2.48,-2.32) {解题图4a：$-\tan\theta=A\sin\theta+B$ 的图解};
\end{tikzpicture}

  \caption{依据答案解题图4a重绘的正切图解。}\label{fig:cpho39q04tan}
\end{figure}
\begin{figure}[htbp]\centering
  \begin{tikzpicture}[>=Stealth,x=1.32cm,y=1.04cm]
  \draw[->] (-0.28,0) -- (5.1,0) node[right] {$\theta$};
  \draw[->] (0,-2.35) -- (0,2.45) node[above] {$y$};
  \foreach \x/\lab in {1.57/{\pi/2},3.14/{\pi},4.71/{3\pi/2}} {
    \draw[densely dashed,gray] (\x,-2.25) -- (\x,2.3);
    \node[below] at (\x,0) {$\lab$};
  }
  \begin{scope}\clip (0,-2.25) rectangle (5.0,2.3);
    \draw[thick,domain=0.07:1.51,samples=100] plot (\x,{-cot(\x r)});
    \draw[thick,domain=1.58:3.06,samples=100] plot (\x,{-cot(\x r)});
    \draw[thick,domain=3.22:4.66,samples=100] plot (\x,{-cot(\x r)});
    \draw[thick,domain=4.73:5.0,samples=50] plot (\x,{-cot(\x r)});
    \draw[thick,blue!65!black,domain=0:5.0,samples=150] plot (\x,{0.88*cos(\x r)+0.72});
  \end{scope}
  \draw[densely dotted] (0,1.60) -- (5.0,1.60) node[right] {$A'+B'$};
  \draw[densely dotted] (0,-0.16) -- (5.0,-0.16) node[right] {$B'-A'$};
  \node[blue!65!black,above] at (3.82,0.16) {$A'\cos\theta+B'$};
  \node[above] at (4.04,1.82) {$-\cot\theta$};
  \node[below=0.68cm] at (2.48,-2.32) {解题图4a$'$：$-\cot\theta=A'\cos\theta+B'$ 的图解};
\end{tikzpicture}

  \caption{依据答案解题图4a$'$重绘的余切图解。}\label{fig:cpho39q04cot}
\end{figure}
由曲线交点可知：第I象限无平衡位形；第II、IV象限各恰有一个。第III象限可能有 $0$、$1$ 或 $2$ 个。临界时两曲线相切，故
\[
A^{2/3}=B^{2/3}+1,
\qquad
\boxed{\left(\frac{4a}{3}\right)^{2/3}-R^{2/3}=\left(\frac{g}{\omega^2}\right)^{2/3}.}
\]
左端小于右端时第III象限无平衡位形；等号时有一个临界位形；左端大于右端时有两个。

\paragraph{（4）受力检验。}
图\ref{fig:cpho39q04forces}按答案解题图4b重绘。重力 $mg$ 作用于杆的中点；离心力是分布力，图中以对 $P$ 点力矩等效的合力 $F_c$ 表示。各象限的力矩方向与上面的交点结论相符。
\begin{figure}[htbp]\centering
  \begin{tikzpicture}[>=Stealth,scale=1.15]
  \begin{scope}
    \draw[densely dashed] (0,-1.05) -- (0,1.55);
    \node[left] at (0,0) {$O$};
    \draw[thick] (0,0) -- (1.55,0) node[midway,above] {$R$};
    \coordinate (P) at (1.55,0); \node[below] at (P) {$P$};
    \draw[very thick,fill=gray!16] (P) -- (2.95,1.02);
    \node[above] at (2.45,0.82) {$2a$}; \node at (1.9,0.48) {$\theta$};
    \draw[->,thick] (2.28,0.54) -- (2.28,-0.62) node[right] {$mg$};
    \draw[->,thick] (2.52,0.72) -- (3.55,0.72) node[right] {$F_c$};
    \node[below,text width=4.4cm,align=center] at (1.65,-1.25) {第I象限：不可能有平衡位形};
  \end{scope}
  \begin{scope}[xshift=7.3cm]
    \draw[densely dashed] (0,-1.05) -- (0,1.55);
    \node[left] at (0,0) {$O$};
    \draw[thick] (0,0) -- (1.55,0) node[midway,above] {$R$};
    \coordinate (P) at (1.55,0); \node[above] at (P) {$P$};
    \draw[very thick,fill=gray!16] (P) -- (2.95,-1.02);
    \node[above] at (2.42,-0.55) {$2a$}; \node at (1.9,0.48) {$\theta$};
    \draw[->,thick] (2.28,-0.54) -- (2.28,-1.7) node[left] {$mg$};
    \draw[->,thick] (2.52,-0.72) -- (3.55,-0.72) node[right] {$F_c$};
    \node[below,text width=4.4cm,align=center] at (1.65,-1.83) {第II象限：可有平衡位形};
  \end{scope}
\end{tikzpicture}

  \caption{依据答案解题图4b重绘的第I、II象限受力及平衡可能性。}\label{fig:cpho39q04forcesupper}
\end{figure}
\begin{figure}[htbp]\centering
  \begin{tikzpicture}[>=Stealth,scale=1.05]
  \begin{scope}
    \draw[densely dashed] (0,-0.98) -- (0,1.5);
    \node[left] at (0,0) {$O$}; \draw[thick] (0,0) -- (1.45,0) node[midway,above] {$R$};
    \coordinate (P) at (1.45,0); \node[above] at (P) {$P$}; \draw[very thick,fill=gray!16] (P) -- (0.12,-0.95);
    \node[below] at (0.58,-0.72) {$2a$}; \node at (1.76,0.47) {$\theta$};
    \draw[->,thick] (0.75,-0.5) -- (0.75,-1.48) node[right] {$mg$}; \draw[->,thick] (0.38,-0.7) -- (1.34,-0.7) node[right] {$F_c$};
    \node[below,text width=5.0cm,align=center,font=\small] at (1.45,-1.72) {第III象限：\(\left(\frac{4a}{3}\right)^{2/3}-R^{2/3}<\left(\frac{g}{\omega^2}\right)^{2/3}\) 时无平衡};
  \end{scope}
  \begin{scope}[xshift=7.2cm]
    \draw[densely dashed] (0,-0.98) -- (0,1.5);
    \node[left] at (0,0) {$O$}; \draw[thick] (0,0) -- (1.45,0) node[midway,above] {$R$};
    \coordinate (P) at (1.45,0); \node[below] at (P) {$P$}; \draw[very thick,fill=gray!16] (P) -- (0.12,0.95);
    \node[above] at (0.62,0.68) {$2a$}; \node at (1.77,0.47) {$\theta$};
    \draw[->,thick] (0.75,0.5) -- (0.75,-0.48) node[right] {$mg$}; \draw[->,thick] (0.38,0.7) -- (-0.58,0.7) node[left] {$F_c$};
    \node[below,text width=5.0cm,align=center,font=\small] at (1.45,-1.25) {第III象限：\(\left(\frac{4a}{3}\right)^{2/3}-R^{2/3}\ge\left(\frac{g}{\omega^2}\right)^{2/3}\) 时可有平衡};
  \end{scope}
\end{tikzpicture}

  \caption{依据答案解题图4b重绘的第III象限受力及平衡可能性。}\label{fig:cpho39q04forcesiii}
\end{figure}
\begin{figure}[htbp]\centering
  \begin{tikzpicture}[>=Stealth,scale=1.25]
  \draw[densely dashed] (0,-1.05) -- (0,1.65);
  \node[left] at (0,0) {$O$}; \draw[thick] (0,0) -- (1.65,0) node[midway,below] {$R$};
  \coordinate (P) at (1.65,0); \node[below right] at (P) {$P$}; \draw[very thick,fill=gray!16] (P) -- (0.15,1.08);
  \node[above] at (0.73,0.78) {$2a$}; \node at (2.02,-0.52) {$\theta$};
  \draw[->,thick] (0.85,0.57) -- (0.85,-0.58) node[right] {$mg$}; \draw[->,thick] (0.45,0.78) -- (1.55,0.78) node[right] {$F_c$};
  \node[below,text width=5.0cm,align=center] at (1.45,-0.92) {第IV象限：可有平衡位形};
\end{tikzpicture}

  \caption{依据答案解题图4b重绘的第IV象限受力及平衡可能性。}\label{fig:cpho39q04forces}
\end{figure}

\paragraph{（5）稳定性。}
在平衡位置 $\theta_E$ 附近令 $\theta=\theta_E+\Delta\theta$。若微小偏移后合力矩为恢复力矩，则平衡稳定。答案的微扰法和有效势能二阶导数法给出一致判断：第II象限位形稳定，第IV象限位形不稳定；第III象限有两个位形时，较大的根稳定、较小的根不稳定；临界相切时 $\theta_E=\theta_0$ 是势能拐点，亦不稳定。

\subsection{转写依据}
题干依据 \texttt{CPHO39-Q04.md} 第2至3页，解答依据同文件答案第12至19页及 \texttt{CPHO39-Q04\_formula\_check\_v1.csv}。原题图4a、答案解题图4a、4a$'$、4b均按对应原页重绘；不嵌入PDF截图。本题五问分值依次为18、3、12、13、14分。
